% =============================================================================
% 1. fejezet — Bevezető
% =============================================================================

\section{Motiváció és problémafelvetés}

A bőrrák napjaink egyik legelterjedtebb daganatos megbetegedése, és bár a
melanoma a bőrdaganatos eseteknek csak kisebb hányadát teszi ki, agresszív
viselkedése és rossz prognózisa miatt a bőrrák okozta halálozás túlnyomó részéért felel.
A betegség kimenetele döntően a felismerés időzítésén múlik: a még lokalizált stádiumban
diagnosztizált melanoma jó eséllyel gyógyítható, késői, áttétes stádiumban
azonban a túlélési esély drasztikusan romlik. A korai felismerés tehát nem csupán
orvosi, hanem népegészségügyi szempontból is kritikus jelentőségű; a téma
részletes epidemiológiai hátterét a \ref{chpt:literature}.~fejezet tárgyalja.

A korai felismerés legnagyobb gyakorlati akadálya a hozzáférés. A hagyományos
diagnosztikai folyamat -- dermatológus által végzett vizuális vizsgálat és
dermatoszkópos értékelés -- magas szintű szakmai tudást és speciális eszközöket
igényel. A világ számos részén, különösen vidéki térségekben és alacsonyabb
jövedelmű országokban, a bőrgyógyászati szakellátáshoz való hozzáférés
korlátozott, ami késedelmes diagnózishoz és rosszabb kimenetelhez vezet; de még
fejlett egészségügyi rendszerekben is jelentős várakozási idővel kell számolni a
szakorvosi vizsgálatig.

A mesterséges intelligencia, azon belül a mély tanulás megjelenése új
lehetőségeket nyitott a bőrelváltozások automatikus osztályozásában. Az úttörő
munkák kimutatták, hogy a megfelelően betanított neurális hálók a
bőrrák-osztályozásban elérhetik a tapasztalt dermatológusok teljesítményét
\cite{esteva2017dermatologist}, és az architektúrák fejlődésével ez a képesség tovább erősödött.
Az okostelefonok elterjedése és fejlődő kamerarendszerei ezt a képességet a beteg
kezébe adhatják: egy ilyen rendszer a felhasználó által készített fotó alapján
jelezheti, ha egy elváltozás szakorvosi vizsgálatot indokol, csökkentve a
késői diagnózisok számát.

\section{A dolgozat célja}

Jelen dolgozat egy bőrelváltozás-elemző
rendszer tervezését és implementációját mutatja be, amely transzfertanulással (transfer learning) 
finomhangolt Vision Transformer (ViT)
\cite{dosovitskiy2021vit} modellre épül -- a modellt az ISIC 2024
versenyadathalmazon \cite{isic2024} tanítottuk bőrelváltozás-osztályozásra. A
dolgozat két, egymást kiegészítő pillérre tagolódik: egy \emph{kutatási} részre,
amely a modell fejlesztését és tudományos kiértékelését tárgyalja, és egy
\emph{mérnöki} részre, amely a betanított modellt egy teljes, használható
alkalmazássá egészíti ki.

A rendszer fő célkitűzései a következők:

\begin{itemize}
    \item Egy bőrelváltozás-osztályozó Vision Transformer modell betanítása az
    ISIC 2024 adathalmazon, a modell validációja, statisztikailag megalapozott
    kiértékelése és viszonyítási alapokhoz mérése.
    \item Egy Android mobilalkalmazás fejlesztése, amely lehetővé teszi a
    felhasználó számára, hogy fényképet készítsen vagy töltsön fel egy
    bőrelváltozásról, és azonnali, AI-alapú elemzést kapjon.
    \item Egy webalkalmazás fejlesztése, amely lehetővé teszi a
    felhasználó számára, hogy fényképet készítsen vagy töltsön fel egy
    bőrelváltozásról, és azonnali, AI-alapú elemzést kapjon.
    \item A feltöltött kép minőségének és relevanciájának automatikus
    ellenőrzése egy nagy nyelvi modell (Gemini) segítségével, kiszűrve az
    értékelhetetlen vagy nem megfelelő felvételeket.
    \item A bőrelváltozás osztályozása \textit{benignus/jóindulatú} vagy \textit{malignus/rosszindulatú}
    kategóriába a ViT-modellel, konfidencia-értékkel kiegészítve.
    \item Természetes nyelvű magyarázat és orvosi szempontú visszajelzés
    generálása a klasszifikációs eredmény alapján, az ABCDE-szabály\footnote{Aszimmetria, határvonal, szín, átmérő, evolúció -- a melanoma
    gyanújelei.} figyelembevételével.
    \item A felhasználói interakciók és a feltöltött képek biztonságos, a
    bejelentkező felhasználóhoz kötött tárolása felhőalapú infrastruktúrán.
\end{itemize}

Fontos hangsúlyozni, hogy a rendszer nem helyettesíti a szakorvosi diagnózist:
kizárólag tájékoztató jellegű szűrőeszközként funkcionál, amely segíti a
felhasználót abban, hogy időben döntsön a dermatológiai konzultáció
szükségességéről.

\section{A dolgozat felépítése}

A dolgozat hat fejezetre tagolódik. A jelen \textbf{Bevezető} után a
\textbf{Szakirodalmi áttekintés} (\ref{chpt:literature}.~fejezet) a
bőrelváltozások gépi tanulással történő osztályozásának előzményeit tekinti át --
a klinikai diagnosztikai gyakorlattól a klasszikus gépi tanuláson és a
konvolúciós hálókon át a Vision Transformer és a multimodális nagy nyelvi
modellek megjelenéséig --, valamint ismerteti az ISIC adathalmazok fejlődését.
Az \textbf{Elméleti háttér} (\ref{chpt:theory}.~fejezet) a rendszer mögötti
kulcstechnológiák elméletét építi fel: a mesterséges neurontól és a hálózatok
tanításától a konvolúciós hálókon, a figyelmi mechanizmuson és a Vision
Transformeren át az osztály-egyensúlytalanság kezeléséig, a kiértékelési
metrikákig és a kommunikációs technológiákig. A \textbf{modell fejlesztése és
kiértékelése} (\ref{chpt:research}.~fejezet) az adathalmaz előkészítését, a ViT modell
betanítását, a viszonyítási alapokat (ResNet-18, EVA02, Gemini), valamint a
kvantitatív eredményeket és a statisztikai elemzést tárgyalja. A \textbf{rendszer
tervezése és megvalósítása} (\ref{chpt:engineering}.~fejezet) az osztályozó modellre épülő rendszer 
megvalósítását mutatja be: a követelményeket, a rendszerarchitektúrát, a
modellkiszolgálást, a backendet, a mobil- és webalkalmazást, valamint a hitelesítést
és a felhőalapú adattárolást.
Végül az \textbf{Összefoglalás és jövőbeli munka} (\ref{chpt:summary}.~fejezet) 
a következtetéseket és a lehetséges továbbfejlesztési irányokat foglalja össze.
